% Schémas TikZ du chapitre « Compression sans perte, à la vitesse du calcul ».
% Formats lus dans le codeur de l'étude 03 (outils/03-entropie/rans/src/lib.rs, decodeur.wgsl) et
% dans fondations/03-entropie-derivations.md, § 1. Les valeurs d'exemple sont dites telles.

\tikzset{
  entmot/.style={draw=black!55, line width=0.3pt, minimum width=1.55cm, minimum height=0.5cm,
    inner sep=0pt, font=\sffamily\scriptsize, fill=#1},
  entmot/.default=Code!14,
  entflux/.style={draw=black!55, line width=0.3pt, minimum width=1.1cm, minimum height=0.5cm,
    inner sep=0pt, font=\sffamily\scriptsize, fill=#1},
  entflux/.default=Expo!12,
  entetage/.style={draw=#1, fill=#1!12, rounded corners=2pt, line width=0.6pt, minimum height=1.05cm,
    text width=2.05cm, align=center, font=\sffamily\scriptsize, inner sep=2pt},
}

% \RangeeBitsE{y}{nom}{liste « longueur/couleur/étiquette »} : rangée de bits étroite (32 bits).
\newcommand{\RangeeBitsE}[3]{%
  \node[anchor=east, font=\sffamily\small\bfseries] at (-0.2,#1) {#2};
  \def\curseur{0}%
  \foreach \lon/\coul/\lab in {#3}{%
    \foreach \b in {1,...,\lon}{%
      \pgfmathsetmacro{\xb}{\curseur + \b - 1}%
      \node[draw=black!55, line width=0.3pt, minimum width=0.34cm, minimum height=0.46cm,
        inner sep=0pt, anchor=west, fill=\coul!28] at (\xb*0.34, #1) {};
    }%
    \pgfmathsetmacro{\xm}{(\curseur + \lon/2)*0.34}%
    \node[font=\sffamily\scriptsize, align=center, below] at (\xm, #1-0.25) {\lab};
    \pgfmathsetmacro{\nouveau}{\curseur + \lon}%
    \xdef\curseur{\nouveau}%
  }%
}

% La tuile du codeur rANS entrelacé : 32 voies × N macro-pas, et ce qui est stocké avec elle.
\newcommand{\SchemaTuileRans}{%
\begin{tikzpicture}[font=\sffamily\scriptsize]
  \foreach \c/\v in {0/0, 1/1, 2/2, 3/3, 5/31}{
    \node[font=\sffamily\scriptsize\bfseries] at (\c*1.6, 0.55) {\EF{lane}{voie} \v};
  }
  \node at (4*1.6, 0.55) {$\cdots$};
  \foreach \r/\m in {0/0, 1/1, 2/2, 4/N-1}{
    \node[anchor=east] at (-0.95, -\r*0.62) {$m = \m$};
    \foreach \c/\v in {0/0, 1/1, 2/2, 3/3, 5/31}{
      \node[entmot] at (\c*1.6, -\r*0.62) {\EF{word}{mot} $32m{+}\v$};
    }
    \node at (4*1.6, -\r*0.62) {$\cdots$};
  }
  \foreach \c in {0,1,2,3,5}{ \node at (\c*1.6, -3*0.62) {$\vdots$}; }
  \draw[decorate, decoration={brace, amplitude=4pt, mirror}, Accent] (-0.78,-2.85) -- (8.78,-2.85)
    node[midway, below=5pt, font=\sffamily\scriptsize, Accent]
    {\EF{each row, one macro-step: 32 words of 32 bits = 128 contiguous output bytes}{chaque rangée, un macro-pas : 32 mots de 32 bits = 128 octets de sortie contigus}};
  % Ce qui est stocké avec la tuile.
  \begin{scope}[xshift=10.0cm]
    \node[entflux=Echelle!20, minimum width=4.6cm, anchor=west] (h) at (0,0.1) {\EF{header: stream offset (4 B)}{en-tête : décalage du flux (4 o)}};
    \node[entflux=Index!16, minimum width=4.6cm, anchor=west] (e) at (0,-0.55) {\EF{initial states $32 \times 32$ bits (128 B)}{états initiaux $32 \times 32$ bits (128 o)}};
    \node[entflux, minimum width=4.6cm, anchor=west, minimum height=1.2cm, text width=4.3cm, align=center] (f) at (0,-1.55)
      {\EF{shared stream: 16-bit words, read in lane order at each sub-step (\cref{fig:flux-partage})}{flux partagé : mots de 16 bits, lus dans l'ordre des voies à chaque sous-pas (\cref{fig:flux-partage})}};
    \node[anchor=north west, text width=4.6cm, font=\sffamily\scriptsize, black!70] at (0,-2.3)
      {\EF{$N = 256$: $132$ bytes for $32\,768$ source bytes, i.e.\ \num{0.403}\,\% of headers}{$N = 256$ : $132$ octets pour $32\,768$ octets de source, soit \num{0.403}\,\% d'en-têtes}};
  \end{scope}
\end{tikzpicture}}

% L'entrée de 32 bits de la table de décodage (4 096 fentes) et le mot des grandeurs E2M1.
\newcommand{\SchemaEntreeTable}{%
\begin{tikzpicture}
  % Numéros des bits aux frontières des champs de chaque rangée (bit 31 à gauche).
  \foreach \b in {31, 17, 16, 4, 3, 0}{
    \node[font=\sffamily\tiny, black!60] at ({(31-\b)*0.34+0.17}, 0.4) {\b};
  }
  \foreach \b in {31, 20, 19, 8, 7, 0}{
    \node[font=\sffamily\tiny, black!60] at ({(31-\b)*0.34+0.17}, -0.95) {\b};
  }
  \RangeeBitsE{0}{v1--v3}{15/Echelle/{\EF{bias (15)}{biais (15)}}, 13/Index/{$f$ (13)}, 4/Code/{\EF{sym.}{sym.}}}
  \RangeeBitsE{-1.35}{v4}{12/Echelle/{\EF{bias (12)}{biais (12)}}, 12/Index/{$f-1$ (12)}, 8/Code/{\EF{symbol (8)}{symbole (8)}}}
  \node[anchor=west, font=\sffamily\scriptsize, text width=4.3cm] at (11.25,-0.65)
    {\EF{bias $=$ slot $-$ cumul$(s)$; the state update $x \leftarrow f\lfloor x/4096\rfloor + \text{bias}$ never reads the symbol field (\cref{lem:etat-table}).}%
        {biais $=$ fente $-$ cumul$(s)$ ; la mise à jour $x \leftarrow f\lfloor x/4096\rfloor + \text{biais}$ ne lit jamais le champ symbole (\cref{lem:etat-table}).}};
  % Le mot des grandeurs doublées de E2M1.
  \def\curseur{0}
  \foreach \g/\c in {12/7, 8/6, 6/5, 4/4, 3/3, 2/2, 1/1, 0/0}{
    \foreach \b in {1,...,4}{
      \pgfmathsetmacro{\xb}{\curseur + \b - 1}
      \node[draw=black!55, line width=0.3pt, minimum width=0.34cm, minimum height=0.46cm, inner sep=0pt,
        anchor=west, fill=Mant!22] at (\xb*0.34, -2.9) {};
    }
    \pgfmathsetmacro{\xm}{(\curseur + 2)*0.34}
    \node[font=\sffamily\scriptsize, fill=Mant!22, inner sep=1.2pt] at (\xm, -2.9) {\g};
    \node[font=\sffamily\tiny, below, black!65] at (\xm, -3.15) {$c\,\&\,7 = \c$};
    \pgfmathsetmacro{\nouveau}{\curseur + 4}
    \xdef\curseur{\nouveau}
  }
  \node[anchor=east, font=\sffamily\small\bfseries] at (-0.2,-2.9) {\texttt{0xC8643210}};
  \node[anchor=west, font=\sffamily\scriptsize, text width=4.3cm] at (11.25,-2.95)
    {\EF{the eight doubled E2M1 magnitudes, one nibble each: $g = (\texttt{0xC8643210} \gg 4(c\,\&\,7))\,\&\,15$, sign from bit 3; no table in memory}%
        {les huit grandeurs E2M1 doublées, un demi-octet chacune : $g = (\texttt{0xC8643210} \gg 4(c\,\&\,7))\,\&\,15$, signe au bit 3 ; aucune table en mémoire}};
\end{tikzpicture}}

% Le flux partagé : vote du sous-groupe et comptage de bits (8 voies montrées sur 32).
\newcommand{\SchemaFluxPartage}{%
\begin{tikzpicture}[font=\sffamily\scriptsize]
  \node[anchor=east, font=\sffamily\scriptsize\bfseries] at (-0.4,0) {\EF{lane $\ell$}{voie $\ell$}};
  \node[anchor=east] at (-0.4,-0.6) {$b_\ell = [x_\ell < L]$};
  \node[anchor=east] at (-0.4,-1.2) {$\sum_{j<\ell} b_j$};
  \node[anchor=east] at (-0.4,-1.8) {\EF{reads}{lit}};
  \foreach \l/\b/\s/\w in {0/0/0/{}, 1/1/0/{$w_{p}$}, 2/1/1/{$w_{p+1}$}, 3/0/2/{}, 4/1/2/{$w_{p+2}$}, 5/0/3/{}, 6/0/3/{}, 7/1/3/{$w_{p+3}$}}{
    \node[draw=black!40, minimum width=0.9cm, minimum height=0.45cm, fill=Fond] (v\l) at (\l*1.15, 0) {\l};
    \ifnum\b=1
      \node[draw=Signe, fill=Signe!15, minimum width=0.9cm, minimum height=0.45cm] at (\l*1.15, -0.6) {1};
      \node[draw=black!30, minimum width=0.9cm, minimum height=0.45cm] at (\l*1.15, -1.2) {\s};
      \node[draw=Expo, fill=Expo!14, minimum width=0.9cm, minimum height=0.45cm] (l\l) at (\l*1.15, -1.8) {\w};
    \else
      \node[draw=black!30, minimum width=0.9cm, minimum height=0.45cm] at (\l*1.15, -0.6) {0};
      \node[draw=black!30, minimum width=0.9cm, minimum height=0.45cm, black!40] at (\l*1.15, -1.2) {\s};
      \node[minimum width=0.9cm, minimum height=0.45cm, black!35] at (\l*1.15, -1.8) {--};
    \fi
  }
  % Le flux, en bas.
  \foreach \i/\lab/\c in {0/{$w_{p-1}$}/black!8, 1/{$w_{p}$}/Expo!14, 2/{$w_{p+1}$}/Expo!14, 3/{$w_{p+2}$}/Expo!14, 4/{$w_{p+3}$}/Expo!14, 5/{$w_{p+4}$}/white, 6/{$\cdots$}/white}{
    \node[entflux=\c] (w\i) at (0.9+\i*1.15, -3.3) {\lab};
  }
  \node[anchor=east] at (-0.4,-3.3) {\EF{shared stream}{flux partagé}};
  \foreach \l/\i in {1/1, 2/2, 4/3, 7/4}{ \draw[-{Latex}, Expo, line width=0.5pt] (w\i.north) -- (l\l.south); }
  \node[anchor=west, text width=4.6cm] at (9.3,-0.9)
    {\EF{One \texttt{ballot} gives the 32 bits $b_\ell$; lane $\ell$ reads at $p + \mathrm{popcount}(\mathrm{ballot} \,\&\, (2^\ell - 1))$, then every lane advances $p$ by $\mathrm{popcount}(\mathrm{ballot})$, here $4$. No prefix sum, no lane-to-lane dependency (\cref{lem:flux-partage}).}%
        {Un \texttt{ballot} donne les 32 bits $b_\ell$ ; la voie $\ell$ lit en $p + \mathrm{popcount}(\mathrm{ballot} \,\&\, (2^\ell - 1))$, puis toutes les voies avancent $p$ de $\mathrm{popcount}(\mathrm{ballot})$, ici $4$. Ni somme préfixe, ni dépendance d'une voie à l'autre (\cref{lem:flux-partage}).}};
\end{tikzpicture}}

% Les échelles E8M0 d'une ligne : médiane, écarts, trois classes (valeurs d'exemple).
\newcommand{\SchemaEchellesLigne}{%
\begin{tikzpicture}[font=\sffamily\scriptsize]
  \node[anchor=east, font=\sffamily\scriptsize\bfseries] at (-0.35,0) {\EF{scale $e$}{échelle $e$}};
  \node[anchor=east] at (-0.35,-0.62) {$\delta = e - \tilde e$};
  \node[anchor=east] at (-0.35,-1.24) {\EF{class}{classe}};
  \foreach \i/\e/\d/\k in {0/121/0/0, 1/121/0/0, 2/122/{+1}/1, 3/121/0/0, 4/120/{-1}/-1, 5/121/0/0, 6/121/0/0, 7/123/{+2}/1, 8/121/0/0, 9/120/{-1}/-1, 10/121/0/0, 11/121/0/0}{
    \ifnum\k=1 \def\coul{SerieB}\def\cl{$>$}\fi
    \ifnum\k=0 \def\coul{black}\def\cl{$=$}\fi
    \ifnum\k=-1 \def\coul{SerieA}\def\cl{$<$}\fi
    \node[draw=black!50, fill=Echelle!22, minimum width=0.95cm, minimum height=0.45cm] at (\i*1.02, 0) {\e};
    \node[draw=black!30, minimum width=0.95cm, minimum height=0.45cm] at (\i*1.02, -0.62) {\d};
    \node[draw=\coul, fill=\coul!14, minimum width=0.95cm, minimum height=0.45cm] at (\i*1.02, -1.24) {\cl};
  }
  \node at (12*1.02, 0) {$\cdots$};
  \draw[decorate, decoration={brace, amplitude=4pt, mirror}, black!60] (-0.48,-1.6) -- (11.7,-1.6)
    node[midway, below=5pt, text width=11cm, align=center]
    {\EF{a row of gpt-oss experts: 2\,880 weights = 90 blocks of 32; its median $\tilde e$ (here 121) is stored as is, 8 bits per row; $\delta$ is coded by rANS, \num{1.07} bits per block instead of 8, and $e = \tilde e + \delta$ reads no other block}%
        {une ligne d'expert de gpt-oss : 2\,880 poids = 90 blocs de 32 ; sa médiane $\tilde e$ (ici 121) est stockée en clair, 8 bits par ligne ; $\delta$ est codé par rANS, \num{1.07} bit par bloc au lieu de 8, et $e = \tilde e + \delta$ ne lit aucun autre bloc}};
  \node[draw=SerieA, fill=SerieA!12, rounded corners=1.5pt, text width=3.3cm, align=center] at (1.6,-3.15)
    {\EF{stream ``$<$'', table $T_{<}$}{flux « $<$ », table $T_{<}$}\\[1pt]\num{12.9}\,\% \EF{of blocks}{des blocs}};
  \node[draw=black!60, fill=black!8, rounded corners=1.5pt, text width=3.3cm, align=center] at (5.6,-3.15)
    {\EF{stream ``$=$'', table $T_{=}$}{flux « $=$ », table $T_{=}$}\\[1pt]\num{74.8}\,\%};
  \node[draw=SerieB, fill=SerieB!12, rounded corners=1.5pt, text width=3.3cm, align=center] at (9.6,-3.15)
    {\EF{stream ``$>$'', table $T_{>}$}{flux « $>$ », table $T_{>}$}\\[1pt]\num{12.3}\,\%};
\end{tikzpicture}}

% Où décoder : à l'arrivée du tronçon, une fois, jamais dans le produit.
\newcommand{\SchemaOuDecoder}{%
\begin{tikzpicture}[font=\sffamily\scriptsize, node distance=0.42cm]
  \node[entetage=Disque] (d) {\textbf{\EF{disk}{disque}}\\\EF{coded, $r$}{codé, $r$}\\\num{2.82}\,\EF{GB/s}{Go/s}};
  \node[entetage=Ram, right=of d] (r) {\textbf{\EF{pinned RAM}{RAM épinglée}}\\\EF{coded}{codé}};
  \node[entetage=Pcie, right=of r] (p) {\textbf{PCIe DMA}\\\EF{coded}{codé}\\\num{12.2}\,\EF{GB/s}{Go/s}};
  \node[entetage=Calcul, right=of p] (k) {\textbf{\EF{decode once}{décoder une fois}}\\rANS v4\\\num{70}\,\EF{GB/s}{Go/s}};
  \node[entetage=Vram, right=of k] (v) {\textbf{\EF{VRAM cache}{cache VRAM}}\\\EF{decoded}{décodé}};
  \node[entetage=Calcul, right=of v] (g) {\textbf{\EF{raw GEMV}{produit brut}}\\\num{209}\,\EF{GB/s}{Go/s}};
  \foreach \a/\b in {d/r, r/p, p/k, k/v, v/g}{ \draw[-{Latex}, line width=0.7pt, black!60] (\a) -- (\b); }
  \node[entetage=Signe, below=0.9cm of v, text width=4.2cm] (f) {\textbf{\EF{fused decode + GEMV}{décodage fondu au produit}}\\\num{57}\,\EF{GB/s, 3.7 times slower, no faster than two passes}{Go/s : 3,7 fois plus lent, pas plus rapide que deux passes}};
  \draw[-{Latex}, Signe, dashed, line width=0.6pt] (p.south) |- (f.west);
  \draw[Signe, line width=1.1pt] ($(f.north west)+(0.15,-0.1)$) -- ($(f.south east)+(-0.15,0.1)$);
  \node[anchor=north west, text width=4.9cm, align=left, black!70] at ($(d.south west)+(0,-0.3)$)
    {\EF{The rate $r$ divides the bytes crossing disk, RAM and link; the card pays $\approx 3.6\,\%$ of its time to decode a token read entirely from disk (\num{2.49}\,/\,\num{70}).}%
        {Le taux $r$ divise les octets qui traversent disque, RAM et lien ; la carte paie $\approx 3{,}6\,\%$ de son temps pour décoder un jeton lu entièrement du disque (\num{2.49}\,/\,\num{70}).}};
\end{tikzpicture}}
